\section{Uniform one-boundary coefficient ratios}\label{uniform-one-boundary-coefficient-ratios}

The coefficient problem uses the same harmonic comparison, but not the
higher-order marked expansion. Log-concavity converts the logarithmic
comparison into adjacent ratios in the interior. Near the top degree we
use an exact intersection number; near degree zero we use bracket
difference estimates. We first prove the uniform equivalent and then,
in Appendix~\ref{sec:quantitative-coefficients}, quantify its error.

Put \(N=3g-2\) and retain the defining normalization \[
 V_{g,1}(L)=\sum_{k=0}^{N}\frac{a_{g,k}}{(2k+1)!}L^{2k},
 \qquad b_{g,k}=[\tau_k]_{g,1}=4^k a_{g,k}.
\] For every valid adjacent index \(0\le k<N\), define
\(t_{g,k}\in(0,\pi)\) by \begin{equation}
\alpha_{g,k}:=\frac{2k+1}{2g-1}
 =\frac{t_{g,k}^{2}\sin t_{g,k}}
        {\sin t_{g,k}-t_{g,k}\cos t_{g,k}}.
 \label{eq:C1}
\end{equation} The assertion to prove is \begin{equation}
\sup_{0\le k<3g-2}
 \left|\frac{a_{g,k}/a_{g,k+1}}{4\pi^2/t_{g,k}^{2}}-1\right|
 \longrightarrow0.
 \label{eq:C2}
\end{equation} Equivalently, the comparison for \(b_{g,k}/b_{g,k+1}\) is
\(\pi^2/t_{g,k}^2\). The ratio at \(k=N\) is undefined and is not
included. In the boundary-length normalization of \(V_{g,1}(L)\), the bottom limit is
four, while the top adjacent ratios diverge.

\subsection{Coefficient identities and further inputs}\label{exact-dictionary-and-the-inputs-needed-here}

We first prove the uniform equivalent \eqref{eq:C2}; Appendix~\ref{sec:quantitative-coefficients} supplies the quantitative error in Theorem~\ref{thm:coefficient}. Let \[
 J_k=\int_{\overline{\mathcal M}_{g,1}}
       \kappa_1^{N-k}\psi_1^k,
 \qquad m=N-k,\quad h=2g-1.
\] The exact coefficient identities are \begin{equation}
a_{g,k}=(2k+1)!!\frac{(2\pi^2)^m}{m!}J_k,
 \qquad
 \frac{a_{g,k}}{a_{g,k+1}}
 =\frac{2\pi^2}{(2k+3)(N-k)}\frac{J_k}{J_{k+1}}.
 \label{eq:C3}
\end{equation} Note that \(J_k\) is a mixed \(\kappa_1\)--\(\psi_1\)
intersection number, not a pure \(\psi\) intersection number.

We use the following consequences of Section~\ref{uniform-comparison}.
Define \[
 H(z)=\frac{3(\sin t-t\cos t)}{t^3},\quad t=\sqrt{2z},
 \qquad \widetilde A_m(h)=[z^m]H(z)^{-h},
\] with the removable value \(H(0)=1\), and \begin{equation}
M_{g,k}=4^k\frac{3^h\Gamma(h)}{4^g}
                  (2\pi^2)^m \widetilde A_m(h).
 \label{eq:C4}
\end{equation} Thus \(H=H_0\) is the kernel of Section~\ref{the-positive-compact-family-of-kernels}
at \(Y=0\). The tilde indicates that \(\widetilde A_m(h)\) is a coefficient of the
unscaled kernel \(H_0\); thus \(\widetilde A_m(h)=r_0^{-m}A_m(h)\), with
\(A_m(h)\) from \eqref{eq:H12} at \(Y=0\) and \(r_0=\pi^2/2\).
There is an absolute \(C\) such that, for every sufficiently large
\(g\), uniformly in \(0\le k\le N\), \begin{equation}
\exp[-C\log^2(g+1)]
 \le\frac{b_{g,k}}{M_{g,k}}
 \le\exp(C\sqrt g).
 \label{eq:C5}
\end{equation} This is the case of one marking of
\eqref{eq:H42}--\eqref{eq:H44}. At \(Y=0\), the unscaled amplitude is
\(U^0=b/[4^k(2\pi^2)^m]\); undoing that scaling gives exactly
\eqref{eq:C4}, without an additional genus-dependent factor. In particular \begin{equation}
\sup_{0\le k\le N}\frac1g
       \left|\log\frac{b_{g,k}}{M_{g,k}}\right|\longrightarrow0.
 \label{eq:C6}
\end{equation} The lower estimate in \eqref{eq:C5} will also control the bottom endpoint.

We also use the following three facts, which are independent of the
harmonic comparison.

\begin{enumerate}[label=\arabic*.]
\item
  The cotangent classes are nef on the compactification. Indeed, the
  nef logarithmic relative dualizing class in~\cite[Theorem~0.4]{Keel}
  is the cotangent class at the additional mark under the identification
  of the universal curve with the moduli space with one additional mark.
  The logarithmic
  kappa class is nef as well: for the forgetful universal curve
  \(\pi:\overline{\mathcal M}_{g,n+1}\to\overline{\mathcal M}_{g,n}\),
  one has \(\kappa_1=\pi_*(\psi_{n+1}^2)\)~\cite{AC}. Its intersection against any
  complete curve is the square of the pulled-back nef cotangent class on
  the pulled-back universal surface, and hence is nonnegative. This
  argument can be carried out after a finite cover supporting a family
  and on its projective coarse surface; passage to the coarse space
  multiplies all mixed intersection numbers by the same positive factor
  and does not affect the inequality.
  Thus the nef rational divisor classes on the projective coarse moduli
  space satisfy the Khovanskii--Teissier inequalities~\cite[Section~1.6]{Laz}; see also~\cite[Theorem~2.15]{BFJ}. This gives \begin{equation}
J_k^2\ge J_{k-1}J_{k+1},\qquad 1\le k<N.
  \label{eq:C7}
\end{equation} Thus \(q_k:=J_k/J_{k+1}\) is nondecreasing. No log-concavity of the
  normalized bracket sequence is assumed.
\item
  Bracket monotonicity and the fixed-marking absolute difference estimate of Mirzakhani--Zograf~\cite[Theorem~3.2(ii)]{MZ} give, with \(c_k=b_{g,k}/V_{g,1}\), \begin{equation}
c_0=1,\qquad 0\le c_k-c_{k+1}\le C\frac{k+1}{g},
  \qquad 0\le k<N.
  \label{eq:C8}
\end{equation} For \(k\ge1\), their explicit bound is
  \(b_{g,k}-b_{g,k+1}\le16((k+1)V_{g,0}+kV_{g-1,2})\); the fixed-marking
  volume ratios turn it into \eqref{eq:C8}. The \(k=0\) case is the fixed-degree
  estimate \(1-c_1=O(1/g)\). Summing \eqref{eq:C8} gives \(1-c_k\le Ck(k+1)/g\).
  These are absolute estimates, so they cannot be divided by an
  arbitrarily small \(c_k\).
\item
  The fixed-marking leading volume asymptotic gives \begin{equation}
V_{g,1}\asymp g^{-1/2}\Gamma(2g-1)
                      (4\pi^2)^{2g-2},
  \label{eq:C9}
\end{equation} and the one-point pure-$\psi$ identity is \(J_N=1/(24^g g!)\). Only
  the upper bound in \eqref{eq:C9} is needed below. The exact value of the top
  ratio is derived below from the ordinary DVV and string equations.
\end{enumerate}

\subsection{The saddle equation for the degree}\label{the-positive-saddle-and-its-effective-density}

The reciprocal of \(H\) has strictly positive Taylor coefficients and is
analytic on a neighborhood of \([0,\pi^2/2]\). These facts also follow
directly from \(2zH''+5H'+H=0\): writing \(-H'/H=\sum_{j\ge0}r_jz^j\)
gives \(r_0=1/5\) and \[
 (2j+5)r_j=2\sum_{u+v=j-1}r_ur_v>0\quad(j\ge1).
\] Hence \(-\log H\) and \(H^{-1}\) have positive nonconstant
coefficients. There is no positive real zero of \(H\) through
\(\pi^2/2\), because \(\sin t-t\cos t>0\) for \(0<t\le\pi\); by
Pringsheim's theorem this excludes a smaller complex singularity of
\(H^{-1}\).

For \(0<z\le\pi^2/2\), tilt these coefficients to a probability law with
generating function \(H(z)/H(zw)\). Its mean is \begin{equation}
\mu(z)=-\frac{zH'(z)}{H(z)}
 =\frac12\left(3-
     \frac{t^2\sin t}{\sin t-t\cos t}\right),
 \qquad t=\sqrt{2z}.
 \label{eq:C10}
\end{equation} Its variance is \(z\mu'(z)>0\). Thus the right side of \eqref{eq:C1}, as a
function of \(t\), decreases strictly from three to zero on \((0,\pi)\),
proving existence and uniqueness in \eqref{eq:C1}. Notice the exact dimension
identity \begin{equation}
\frac mh=\frac{3-\alpha_{g,k}}2,
 \qquad \alpha_{g,k}=3-\frac{2m}{h}.
 \label{eq:C11}
\end{equation}

For \(m/h\) in a compact subset of \((0,3/2)\), there is a unique saddle
\(z\) with \(\mu(z)=m/h\). Fourier inversion for the tilted iid sum,
whose laws have uniform exponential moments and positive probabilities
at zero and one, gives \begin{equation}
\log \widetilde A_m(h)=-h\log H(z)-m\log z+O(\log h).
 \label{eq:C12}
\end{equation} The error is uniform on such compact subsets.

\subsection{Interior degrees: exponential comparison and concavity}\label{interior-degrees-exponential-comparison-plus-concavity}

Suppose \(k/g\to\alpha\in(0,3)\) and let \(z=t^2/2\) solve
\(\mu(z)=(3-\alpha)/2\). Interpolate linearly on the degree mesh the
values \begin{equation}
f_g(k/g)=\frac1g\log
               \frac{J_k}{J_Ng^{2(N-k)}}.
 \label{eq:C13}
\end{equation} By \eqref{eq:C7}, these functions are concave. Equations \eqref{eq:C3}--\eqref{eq:C6}, \eqref{eq:C12},
and Stirling's formula give their pointwise limit on \((0,3)\): \begin{equation}
\begin{aligned}
 f(\alpha)={}&-2\log H(z)-(3-\alpha)\log z
 +(3-\alpha)\log(3-\alpha)\\
 &+3\log6-\alpha\log(2\alpha)-2(3-\alpha).
 \end{aligned}
 \label{eq:C14}
\end{equation} Explicitly, the model for the intersection number \(J_k\) is
\(J_k^{\mathrm{mod}}=3^h\Gamma(h)4^{-g}m!\widetilde A_m(h)/(2k+1)!!\). Its quotient
by the top model value is exactly \(m! \widetilde A_m(h)(2N+1)!!/(2k+1)!!\); the
two occurrences of \((N-k)\log g\) in its Stirling expansion cancel the
affine normalization in \eqref{eq:C13}. The logarithmic errors in \eqref{eq:C6}, including
the top one, vanish after division by \(g\).

Differentiation of \eqref{eq:C14} gives \begin{equation}
f'(\alpha)=\log\frac{z}{2\alpha(3-\alpha)}.
 \label{eq:C15}
\end{equation} The terms containing \(dz/d\alpha\) cancel by the saddle equation.
Concavity now upgrades the exponential-scale information to an adjacent
ratio. For fixed small \(\varepsilon>0\), the slope on \([k/g,(k+1)/g]\)
is at most the secant slope on
\([\alpha-2\varepsilon,\alpha-\varepsilon]\) and at least the one on
\([\alpha+\varepsilon,\alpha+2\varepsilon]\), for all sufficiently large
\(g\). These four endpoints are fixed, so pointwise convergence
suffices. First let \(g\) tend to infinity and then let \(\varepsilon\)
decrease to zero. Since \eqref{eq:C14} is differentiable, the neighboring slope
converges to \eqref{eq:C15}. Consequently \begin{equation}
\frac{J_k}{g^2J_{k+1}}\longrightarrow
            \frac{2\alpha(3-\alpha)}z.
 \label{eq:C16}
\end{equation} Using the exact identity \eqref{eq:C3}, \begin{equation}
\frac{a_{g,k}}{a_{g,k+1}}\longrightarrow
      \frac{2\pi^2}{z}=\frac{4\pi^2}{t^2}.
 \label{eq:C17}
\end{equation} The passage from macroscopic secants to adjacent ratios uses the
concavity of $\log J_k$.

\subsection{The top endpoint, including bounded codimension}\label{the-top-endpoint-including-bounded-codimension}

The one-kappa pushforward identity gives
\(J_{N-1}=\langle\tau_2\tau_{3g-3}\rangle_g\). Applying DVV at the
degree-two insertion and then the string equation twice gives \[
 15J_{N-1}=
 \{(6g-3)(6g-5)+12g\}J_N.
\] The nonseparating term is
\((1/2)\langle\tau_0^2\tau_{3g-3}\rangle_{g-1}=12gJ_N\); separating
terms vanish by dimension. At \(g=1\), the same identity follows
directly from \(\langle\tau_0^3\rangle_0=1\). Therefore \begin{equation}
q_{N-1}=\frac{12g^2-12g+5}{5},
 \qquad
 \frac{a_{g,N-1}}{a_{g,N}}
 =\frac{2\pi^2(12g^2-12g+5)}{5(6g-3)}.
 \label{eq:C18}
\end{equation}

Let \(m=N-k\ge1\) and \(m=o(g)\). For each fixed \(\alpha<3\),
eventually \(q_{\lfloor\alpha g\rfloor}\le q_{N-m}\le q_{N-1}\) by \eqref{eq:C7}.
As \(\alpha\) increases to three, the limit on the right side of \eqref{eq:C16}
tends to \(12/5\), because \begin{equation}
\frac{t^2\sin t}{\sin t-t\cos t}
       =3-\frac{t^2}{5}+O(t^4)\quad(t\to0).
 \label{eq:C19}
\end{equation} The exact top value \eqref{eq:C18} gives the same limit. Hence \begin{equation}
\frac{q_{N-m}}{g^2}\longrightarrow\frac{12}{5},
 \qquad
 \frac{a_{g,N-m}}{a_{g,N-m+1}}
       \sim\frac{4\pi^2g}{5m}\quad(m=o(g),\ m\ge1).
 \label{eq:C20}
\end{equation} In particular the statement includes every fixed positive
codimension.
For fixed $m$, this leading term agrees, after the length normalization
above, with the fixed-codimension expansion of
Liu--Xu~\cite[Theorem~1.2]{LX}. The same monotonicity argument also covers
unbounded $m=o(g)$ and joins this endpoint to the interior limit.

\subsection{The lowest degrees}\label{the-bottom-endpoint-without-a-relative-error-assumption}

Near degree zero the interior argument gives no information: the
comparison \eqref{eq:C5} determines \(\log b_{g,k}\) only up to an error
\(o(g)\), whereas for \(k=o(g)\) the quantities to be determined vary on a
smaller scale. We combine two other inputs. For degrees up to a power of \(\log g\), the absolute estimate
\eqref{eq:C8} of Mirzakhani and Zograf gives the ratios directly. For
larger \(k=o(g)\), we argue by contradiction: by the monotonicity of the
ratios \(q_k\), a ratio bounded away from one would persist over a range
of degrees comparable to \(k\), and would force \(b_{g,k}\) to drop below
the lower bound \eqref{eq:C23}.

Set \(z_0=\pi^2/2\), and let \(S_h\) be the sum of \(h\) iid variables
with generating function \(H(z_0)/H(z_0w)\). Since \(H(z_0)=3/\pi^2\),
\eqref{eq:C4} becomes the exact identity \begin{equation}
M_{g,k}=\pi^2\Gamma(h)(4\pi^2)^{h-1}
                     \Pr(S_h=N-k).
 \label{eq:C21}
\end{equation} The summands have mean \(3/2\), nonzero variance, span one, and
exponential moments. Uniform exponential tilting in a fixed neighborhood
of that mean gives, for \(0\le k\le c_0g\) with one fixed small
\(c_0>0\), \begin{equation}
\Pr(S_h=N-k)\ge c g^{-1/2}
            \exp[-C(k+1)^2/g].
 \label{eq:C22}
\end{equation} Indeed the tilted local central limit estimate gives the factor
\(g^{-1/2}\), while the rate function has a quadratic zero at the mean;
the displacement here is exactly \(k+1/2\). Combining \eqref{eq:C5}, \eqref{eq:C9}, and
\eqref{eq:C21}--\eqref{eq:C22} yields \begin{equation}
0\le\log\frac{V_{g,1}}{b_{g,k}}
 \le C\log^2(g+1)+C\frac{(k+1)^2}{g}.
 \label{eq:C23}
\end{equation}

Write \[
 \rho_k=\frac{b_{g,k}}{b_{g,k+1}}\ge1,
 \qquad \omega_k=(2k+3)(N-k),
 \qquad \rho_k=\frac{\pi^2q_k}{2\omega_k}.
\] Suppose \(k=o(g)\) but, along a subsequence, \(\rho_k\ge1+\eta\) for
a fixed \(\eta>0\). First consider \(k>\log^3(g+1)\); then \eqref{eq:C23} is
\(o(k)\) at every index comparable to \(k\). Put
\(\delta=\eta/[2(2+\eta)]\). For \(k\le j\le\lfloor(1+\delta)k\rfloor\),
\begin{equation}
\frac{\omega_j}{\omega_k}\le\frac{2j+3}{2k+3}\le1+\delta,
 \qquad
 \rho_j\ge\rho_k\frac{\omega_k}{\omega_j}
            \ge1+\eta/2.
 \label{eq:C24}
\end{equation} These inequalities follow from the monotonicity of the raw ratios $q_j$.
Multiplying these inequalities forces
\(\log(V_{g,1}/b_{g,\lfloor(1+\delta)k\rfloor+1})\ge c_\eta k\),
contradicting \eqref{eq:C23}.

For the remaining indices \(k\le\log^3(g+1)\), summing the absolute
estimate \eqref{eq:C8} gives \(c_{k+1}=1-o(1)\) uniformly, so division in \eqref{eq:C8} is
legitimate. In fact there is a fixed small \(\delta_0>0\) for which \begin{equation}
\rho_k=1+O((k+1)/g)
       \quad(0\le k\le\delta_0\sqrt g).
 \label{eq:C25}
\end{equation} This proves \begin{equation}
\frac{a_{g,k}}{a_{g,k+1}}\longrightarrow4
                  \quad\text{whenever }k=o(g).
 \label{eq:C26}
\end{equation} No positive lower bound on \(c_k\) was assumed outside the small
square-root window.

\subsection{Combining the degree ranges}\label{uniformity-and-the-corrected-top-density}

\begin{proof}[Proof of the uniform equivalent \eqref{eq:C2}]
Take any sequence of valid indices and a subsequence on which \(k/g\)
converges in \([0,3]\). An interior limit is covered by \eqref{eq:C17}, and a
zero limit by \eqref{eq:C26}, while the comparison in \eqref{eq:C1} also tends to four. If
the limit is three, then \(m=N-k=o(g)\) and \eqref{eq:C20} applies. Equations
\eqref{eq:C11} and \eqref{eq:C19} give \[
 t_{g,k}^2\sim\frac{10m}{2g-1},\qquad
 \frac{4\pi^2}{t_{g,k}^2}\sim\frac{4\pi^2g}{5m}.
\] Thus the true and the proposed ratios are asymptotically equal on that
subsequence too. A sequence violating the supremum statement \eqref{eq:C2} would
contradict one of these three cases.
\end{proof}

The effective density \((2k+1)/(2g-1)\) in \eqref{eq:C1} is essential near the
top: it retains the exact codimension \(m\) in
\(3-\alpha_{g,k}=2m/(2g-1)\). Replacing it by \(k/g\) before taking the
endpoint limit would change the fixed-codimension denominator.
